% TEMPLATE-MILC-v3.tex - plantilla maestra para todas las guias MILC v3
% Ver scripts/generadores/build-guia-milc.py para la lista completa de
% claves esperadas. El script valida que no queden placeholders sin
% reemplazar antes de escribir el archivo final.

\documentclass[11pt,letterpaper]{article}
\usepackage[margin=1.75cm]{geometry}
\usepackage[table]{xcolor}
\usepackage{fontspec}
\usepackage[spanish]{babel}
\usepackage{tikz}
\usepackage[most]{tcolorbox}
\usepackage{tabularx}
\usepackage{array}
\usepackage{fancyhdr}
\usepackage{titlesec}
\usepackage{ragged2e}
\usepackage{enumitem}
\usepackage{microtype}

\setmainfont{Helvetica Neue}
\setsansfont{Helvetica Neue}
\setlength{\parindent}{0pt}
\setlength{\parskip}{6pt}
\hyphenpenalty=200
\exhyphenpenalty=200
\pretolerance=400
\tolerance=2000
\setlength{\emergencystretch}{1em}
\renewcommand{\arraystretch}{1.25}
\setlist{leftmargin=5mm,itemsep=1mm,topsep=1mm}
\newcolumntype{Y}{>{\RaggedRight\arraybackslash}X}

\definecolor{milcMagenta}{HTML}{E6007E}
\definecolor{milcVino}{HTML}{5A0038}
\definecolor{milcTurquesa}{HTML}{0093A5}
\definecolor{milcVerde}{HTML}{58A923}
\definecolor{milcMostaza}{HTML}{E5B400}
\definecolor{milcOcre}{HTML}{B86600}
\definecolor{milcNegro}{HTML}{241816}
\definecolor{milcGris}{HTML}{F7F7F4}
\definecolor{milcLinea}{HTML}{E8E2DD}
\definecolor{milcGradePrimary}{HTML}{84CC16}
\definecolor{milcGradeDark}{HTML}{4D7C0F}
\definecolor{milcGradeSoft}{HTML}{ECFCCB}
\definecolor{milcGradeAccent}{HTML}{CFE7FF}
\definecolor{milcGradeText}{HTML}{FFFFFF}
\color{milcNegro}

\pagestyle{fancy}
\fancyhf{}
\lhead{\small Tecnología e Informática · Grado 7}
\rhead{\small Guía 1 · MILC v3}
\cfoot{\small\thepage}
\renewcommand{\headrulewidth}{0pt}

\newtcolorbox{titlebox}[1]{
  enhanced,colback=#1,colframe=#1,arc=7mm,boxrule=0pt,
  left=7mm,right=7mm,top=3mm,bottom=3mm,
  before skip=8pt,after skip=8pt,
  fontupper=\sffamily\bfseries\Large\color{white}
}

\newtcolorbox{softbox}[3]{
  enhanced,colback=#2,colframe=#1,borderline west={4pt}{0pt}{#1},
  arc=2mm,boxrule=.4pt,left=4mm,right=4mm,top=3mm,bottom=3mm,
  before skip=6pt,after skip=8pt,title={#3},coltitle=white,
  colbacktitle=#1,fonttitle=\sffamily\bfseries,fontupper=\RaggedRight,
  attach boxed title to top left={xshift=3mm,yshift=-2mm},
  boxed title style={arc=2mm, boxrule=0pt}
}

\newtcolorbox{drawbox}[2]{
  enhanced,colback=white,colframe=#1,arc=4mm,boxrule=1pt,
  left=5mm,right=5mm,top=4mm,bottom=4mm,before skip=8pt,after skip=8pt,
  title={#2},coltitle=white,colbacktitle=#1,fonttitle=\sffamily\bfseries,
  fontupper=\RaggedRight,
  attach boxed title to top left={xshift=4mm,yshift=-2mm},
  boxed title style={arc=3mm, boxrule=0pt}
}

\newtcolorbox{infoband}[1]{
  enhanced,colback=#1!8,colframe=#1!50,arc=2mm,boxrule=.5pt,
  left=4mm,right=4mm,top=2mm,bottom=2mm,before skip=4pt,after skip=4pt,
  fontupper=\small\RaggedRight
}

% Puente narrativo entre secciones (contrato editorial Conéctate).
% Da continuidad: "Acabas de X. Ahora vas a Y porque Z."
% Visualmente: banda izquierda en color del grado, texto en itálica pequeña.
\newtcolorbox{puentebox}{
  enhanced,colback=milcGradeSoft,colframe=milcGradeDark,
  borderline west={3pt}{0pt}{milcGradeDark},
  arc=0mm,boxrule=0pt,left=5mm,right=4mm,top=2.5mm,bottom=2.5mm,
  before skip=4pt,after skip=8pt,
  fontupper=\itshape\small\color{milcNegro}\RaggedRight
}
\newcommand{\puente}[1]{%
  \begin{puentebox}%
    \textcolor{milcGradeDark}{\textbf{\textrightarrow}}\hspace{2mm}#1%
  \end{puentebox}%
}

\newcommand{\linead}[1]{\noindent\color{milcLinea}\rule{#1}{0.5pt}\color{milcNegro}}
\newcommand{\respuesta}[1]{\par\vspace{2mm}\foreach \n in {1,...,#1}{\noindent\color{milcLinea}\rule{\linewidth}{0.5pt}\par\vspace{5mm}}\color{milcNegro}}
\newcommand{\checkbox}{\textcolor{milcGradeDark}{\fbox{\phantom{x}}}\hspace{5pt}}

\newcommand{\phasecard}[4]{
\begin{tcolorbox}[enhanced,colback=#3,colframe=#2,arc=3mm,boxrule=.5pt,height=3.05cm,valign=center,halign=center,left=1.5mm,right=1.5mm,top=2mm,bottom=2mm]
{\sffamily\bfseries\fontsize{22}{24}\selectfont\textcolor{#2}{#1}}\par
{\sffamily\bfseries\small #4}
\end{tcolorbox}
}

\begin{document}
\thispagestyle{empty}
\begin{tikzpicture}[remember picture,overlay]
  \fill[white] (current page.south west) rectangle (current page.north east);
  \path[fill=milcGradePrimary,rounded corners=16mm]
    ([xshift=.55cm,yshift=-.55cm]current page.north west) rectangle
    ([xshift=-.55cm,yshift=3.65cm]current page.south east);

  \node[anchor=north west,text=milcGradeText] at ([xshift=2.0cm,yshift=-1.85cm]current page.north west)
    {\sffamily\bfseries\fontsize{14}{16}\selectfont GRADO};
  \node[anchor=north west,text=milcGradeText,opacity=.82] at ([xshift=2.0cm,yshift=-2.75cm]current page.north west)
    {\sffamily\bfseries\fontsize{9}{11}\selectfont SERIE GUÍAS · TECNOLOGÍA E INFORMÁTICA};

  \begin{scope}[shift={([xshift=-3.05cm,yshift=-2.55cm]current page.north east)}]
    \fill[white,opacity=.80] (-.62,-.52) rectangle (.62,.52);
    \draw[milcGradeText,line width=1pt] (-.62,-.52) rectangle (.62,.52);
    \fill[milcGradeDark] (-.42,-.38) rectangle (-.22,.06);
    \fill[milcGradeAccent] (-.10,-.38) rectangle (.10,.24);
    \fill[milcGradePrimary!60!black] (.22,-.38) rectangle (.42,.38);
  \end{scope}

  \node[anchor=west,text=milcGradeText] at ([xshift=1.9cm,yshift=-12.85cm]current page.north west)
    {\sffamily\bfseries\fontsize{124}{124}\selectfont 7};
  \draw[milcGradeText,line width=1.7pt]
    ([xshift=4.52cm,yshift=-11.68cm]current page.north west) circle (.13cm);

  \node[anchor=west,text=milcGradeText,text width=8.7cm,align=left] at ([xshift=10.35cm,yshift=-11.6cm]current page.north west)
    {
     {\sffamily\bfseries\fontsize{19}{22}\selectfont Apertura\\periodo 1 ---\\la minga\\digital}\\[4mm]
     {\sffamily\bfseries\fontsize{23}{26}\selectfont Guía 1-7-TIC}\\[2mm]
     {\sffamily\fontsize{13}{16}\selectfont Período 1 · Trabajo colaborativo en la nube}\\[5mm]
     {\sffamily\bfseries\fontsize{11}{13}\selectfont Institución Educativa Sor María Juliana}\\[1mm]
     {\sffamily\fontsize{10}{12}\selectfont Autor: Dr. Álvaro Cárdenas Orozco}
    };

  \path[fill=milcGradeSoft,rounded corners=7mm]
    ([xshift=1.35cm,yshift=.85cm]current page.south west) rectangle
    ([xshift=-1.35cm,yshift=4.25cm]current page.south east);
  \draw[milcGradeDark,line width=.65pt,rounded corners=7mm]
    ([xshift=1.35cm,yshift=.85cm]current page.south west) rectangle
    ([xshift=-1.35cm,yshift=4.25cm]current page.south east);
  \node[anchor=center,text=milcNegro,text width=17.0cm,align=center] at ([yshift=2.55cm]current page.south)
    {\sffamily\fontsize{10}{12}\selectfont Metodología MILC · Apertura ancestral · Pensamiento computacional · Triángulo Dussel-Estoicismo-Floridi\\
    \textcolor{milcGradeDark}{\bfseries Producto:} Tu \textbf{cuaderno inaugural de grado 7} listo para el periodo 1: portada con nombre + grado + año + materia; \textbf{índice del periodo 1} con las 10 sesiones del periodo; \textbf{compromiso firmado} con 5 acuerdos de trabajo colaborativo (responsabilidad, puntualidad, respeto, escucha, comunicación clara); \textbf{5 preguntas-radar} sobre trabajo colaborativo y la nube; y \textbf{una mini-investigación} en el cuaderno: \emph{``¿qué es una minga y por qué funcionaba en el campo?''} (puedes preguntar a tus papás o abuelos).};
\end{tikzpicture}
\mbox{}
\newpage

\section*{Apertura: Apertura periodo 1 --- vamos a hacer minga digital}

\begin{softbox}{milcGradeDark}{milcGradeSoft}{Saber ancestral}
\textbf{Cuando un campesino del Valle del Cauca necesitaba cosechar 5 hectáreas de café, no lo hacía solo.} En las veredas alrededor de Cartago, antes de las máquinas grandes, existía \textbf{la minga} (o \emph{convite}, o \emph{mano de vuelta}, según la zona). El dueño del cafetal mandaba aviso a sus vecinos un sábado en la mañana. Llegaban \textbf{20 o 30 personas} con sus canastas. Cada uno se ubicaba en un surco. \textbf{Nadie estorbaba al otro.} Todos cosechaban al mismo tiempo, en partes distintas, pero trabajando \emph{en la misma obra}. Las mujeres preparaban almuerzo grande --- sancocho, arroz, plátano --- y al mediodía se sentaban todos juntos a comer. La minga terminaba en 1 día lo que una sola persona habría hecho en un mes. Y al siguiente sábado, los mismos vecinos iban a hacer minga en la finca del vecino que vino a la suya. \textbf{Esa idea --- muchos trabajando juntos en una sola obra, sin pelearse, organizados, con reciprocidad --- es exactamente lo que hoy llamamos ``trabajo colaborativo en la nube''}. Microsoft 365 y OneDrive son la minga digital. Lo nuevo es la herramienta; lo viejo es la sabiduría.
\end{softbox}

\begin{softbox}{milcTurquesa}{milcGris}{Saber contemporáneo}
Este periodo te enseña a \textbf{trabajar con otros en internet}, no por separado. Hasta ahora tus trabajos los hacías solo: abrías tu Word, escribías tu cosa, la entregabas. Pero \emph{en la universidad y en el trabajo se hace al revés}: equipos de personas trabajan al mismo tiempo en el mismo documento, sin chocar. Esto es posible gracias a la \textbf{nube} (\emph{cloud}) y a \textbf{Microsoft 365}: un paquete de aplicaciones (\emph{Word, Excel, PowerPoint, OneNote, Teams, Outlook, OneDrive}) que vive en internet y se comparte entre personas. \textbf{Tu colegio te da una cuenta institucional gratuita}: la usarás todo el año. Al final del periodo 1, sabrás compartir documentos, escribir en grupo, comentar el trabajo de otros, recuperar versiones anteriores, y entregar un mini-proyecto colaborativo con bitácora. Esas habilidades te las exigirán toda la vida.
\end{softbox}

\begin{softbox}{milcMostaza}{milcGris}{Pregunta-puente}
\textit{Si tu profe te pide hacer un trabajo de Sociales con 3 compañeros, y cada uno vive en una casa distinta, ¿\textbf{cómo lo harían}? ¿Cada uno hace su parte y se junta al final? ¿Hay una manera mejor donde los 4 trabajen al mismo tiempo sin verse?}
\end{softbox}

\begin{softbox}{milcVerde}{milcGris}{Saber y saber hacer}
Al terminar la clase: (1) podrás \textbf{identificar} los 5 principios de la minga aplicados al trabajo digital; (2) sabrás \textbf{explicar} qué es Microsoft 365 y la nube; (3) podrás \textbf{aplicar} los 5 acuerdos del trabajo colaborativo; (4) habrás \textbf{creado} tu cuaderno inaugural firmado + investigación sobre la minga.
\end{softbox}

\small
\begin{tabularx}{\textwidth}{>{\bfseries}p{3.0cm}Y}
\rowcolor{milcGradePrimary!12}Autor & Dr. Álvaro Cárdenas Orozco\\
DBA & Utiliza herramientas digitales para trabajar de manera colaborativa con otros (MEN, T\&I 7°)\\
Referentes & Saber ancestral de la minga · Microsoft 365 y OneDrive · Coautoría asincrónica\\
Producto final & Tu \textbf{cuaderno inaugural de grado 7} listo para el periodo 1: portada con nombre + grado + año + materia; \textbf{índice del periodo 1} con las 10 sesiones del periodo; \textbf{compromiso firmado} con 5 acuerdos de trabajo colaborativo (responsabilidad, puntualidad, respeto, escucha, comunicación clara); \textbf{5 preguntas-radar} sobre trabajo colaborativo y la nube; y \textbf{una mini-investigación} en el cuaderno: \emph{``¿qué es una minga y por qué funcionaba en el campo?''} (puedes preguntar a tus papás o abuelos).\\
\end{tabularx}
\normalsize

\puente{Hoy haces 4 cosas: descubres qué era la minga, conoces la ruta del periodo, firmas el compromiso colaborativo, y dejas el cuaderno listo.}

\begin{titlebox}{milcGradeDark}
Ruta MILC: así vamos a aprender
\end{titlebox}
\begin{tabularx}{\textwidth}{>{\centering\arraybackslash}X>{\centering\arraybackslash}X>{\centering\arraybackslash}X>{\centering\arraybackslash}X}
\phasecard{1}{milcVerde}{milcGris}{Escuta\\[-1mm]\small Observo, escucho y pregunto.} &
\phasecard{2}{milcTurquesa}{milcGris}{Sistematización\\[-1mm]\small Pensamiento computacional.} &
\phasecard{3}{milcMagenta}{milcGris}{Praxis\\[-1mm]\small Construyo y aplico.} &
\phasecard{4}{milcVino}{milcGris}{Evaluación\\[-1mm]\small Reflexión liberadora.}
\end{tabularx}


\puente{Empezamos imaginando una minga real con tus vecinos.}

\begin{titlebox}{milcVerde}
Fase 1 · Escuta
\end{titlebox}

\begin{softbox}{milcVerde}{milcGris}{Escena para escuchar}
\textbf{Actividad 1 · IDENTIFICA --- Imagina la minga} (12 min · individual). Cierra los ojos y piensa en una minga real. \textbf{En el cuaderno responde:} \emph{(1) ¿Cuántas personas hay en la finca un día de minga?} \emph{(2) ¿Cómo se reparten el trabajo --- todos en lo mismo o cada uno en algo distinto?} \emph{(3) ¿Quién organiza?} \emph{(4) ¿Qué pasaría si 5 personas tratan de cosechar el mismo árbol al tiempo?} \emph{(5) ¿Cómo deciden cuándo descansar?} \emph{(6) ¿Qué pasa si alguien no llega o llega tarde?} \emph{(7) ¿Quién prepara el almuerzo?} \emph{(8) ¿Qué pasa al final del día?}. Respuesta con tu intuición. Si tienes abuelos o tíos que hayan vivido en el campo, te pueden contar más después. Esta intuición es la base del trabajo colaborativo digital de hoy.
\end{softbox}

\checkbox Cerré los ojos e imaginé la minga.\\
\checkbox Respondí las 8 preguntas con mi intuición.\\
\checkbox Pensé sin copiar al compañero.

\begin{infoband}{milcVerde}
\textbf{Lo que va al cuaderno (Actividad 1 · IDENTIFICA).} Título: «Actividad 1 --- Imagino la minga». \textbf{Formato:} 8 respuestas cortas. \textbf{Extensión:} 8 líneas. \textbf{Sabes que terminaste cuando} tus respuestas reflejan cómo funcionaría una minga real, no una versión imaginada.
\end{infoband}

\puente{Después aprendes los 5 principios de la minga y los aplicamos al trabajo digital.}

\begin{titlebox}{milcTurquesa}
Fase 2 · Sistematización con pensamiento computacional
\end{titlebox}

\begin{softbox}{milcTurquesa}{milcGris}{Cuatro pilares aplicados al tema}
De tus respuestas salen los \textbf{5 principios de la minga} que aplicamos al trabajo colaborativo digital. Estos principios son antiguos pero universales: funcionaban en 1920 en el café, funcionan en 2026 en Microsoft 365.
\end{softbox}

\small
\begin{tabularx}{\textwidth}{>{\bfseries\RaggedRight\arraybackslash}p{3.6cm}Y}
\rowcolor{milcTurquesa!90}\color{white}Pilar & \color{white}\bfseries Aplicación al tema\\
1. Descomposición & \textbf{Principio 1 --- División sin choque.} En la minga cada quien tomaba un surco distinto. Nadie cosechaba el mismo árbol que otro. \textbf{En coautoría digital es igual:} cada compañero escribe en su párrafo o sección, sin pisarse. Microsoft 365 te muestra dónde está escribiendo cada uno (su cursor aparece con su nombre). \textbf{Principio 2 --- Organización clara.} Alguien tenía que decir \emph{``ustedes 5 al surco de la izquierda, los otros 5 al de la derecha''}. Sin organización inicial, 30 personas se vuelven caos. \textbf{En digital:} antes de empezar a escribir en grupo, se acuerdan los roles (quién escribe qué, qué tiempos, qué entrega cada uno).\\
2. Reconocimiento de patrones & \textbf{Principio 3 --- Reciprocidad.} En la minga, el que recibía ayuda tenía que dar ayuda después. Nadie venía a una minga ajena sin pensar que después le tocaría devolver el favor. \textbf{En digital:} cuando 3 compañeros te ayudan a hacer tu trabajo, después tú vas y los ayudas en los suyos. No es competencia, es reciprocidad. \textbf{Principio 4 --- Comunicación constante.} Durante la minga la gente se hablaba: \emph{``Cuidado con esa rama''}, \emph{``Pase la canasta''}, \emph{``Vamos a almorzar''}. \textbf{En digital:} comentarios en el documento, mensajes en Teams, correos en Outlook. Si nadie comunica, el trabajo se rompe.\\
3. Abstracción & \textbf{Principio 5 --- Cierre y celebración.} La minga terminaba con almuerzo grande y conversación. Era importante: \emph{cerraba el día, agradecía el esfuerzo, sembraba el siguiente encuentro}. \textbf{En digital:} cuando un proyecto termina, hay que cerrarlo bien --- una bitácora, una entrega, un agradecimiento al equipo. Sin ese cierre, los proyectos se sienten incompletos. \textbf{Los 5 acuerdos del trabajo colaborativo (que firmarás).} (1) Responsabilidad: si me asignan algo, lo hago. (2) Puntualidad: si quedamos para una hora, llego. (3) Respeto: si discrepo, lo digo con cuidado. (4) Escucha: pongo atención a lo que dicen los demás. (5) Comunicación clara: digo lo que pienso sin esconderlo.\\
4. Algoritmo conceptual & \textbf{La ruta del periodo 1 --- las 10 sesiones.} \emph{(S1) Hoy: Apertura --- la minga digital.} \emph{(S2) Microsoft 365 --- qué es la nube.} \emph{(S3) OneDrive --- mi nube personal.} \emph{(S4) Word, Excel y PowerPoint en línea.} \emph{(S5) Compartir y permisos.} \emph{(S6) Coautoría asincrónica.} \emph{(S7) Comentarios y sugerencias.} \emph{(S8) Versiones e historial.} \emph{(S9) Outlook y Teams.} \emph{(S10) Cosecha --- mi primer proyecto colaborativo}. Al final del periodo, tu equipo de 3 entrega un proyecto real hecho en la nube.\\
\end{tabularx}
\normalsize

\begin{drawbox}{milcTurquesa}{5 principios de la minga + ruta P1}
\textbf{5 principios:} \\[1mm]
\textbf{1.} División sin choque. \\[1mm]
\textbf{2.} Organización clara. \\[1mm]
\textbf{3.} Reciprocidad. \\[1mm]
\textbf{4.} Comunicación constante. \\[1mm]
\textbf{5.} Cierre y celebración. \\[1mm]
\textbf{Ruta P1:} apertura → Microsoft 365 → OneDrive → Word/Excel/PPT → compartir → coautoría → comentarios → versiones → Outlook/Teams → cosecha.
\end{drawbox}

\begin{softbox}{milcMostaza}{milcGris}{Errores comunes y su corrección}
\textbf{Errores frecuentes al iniciar el trabajo colaborativo.} (1) \emph{``Yo solo, mejor''} --- Falso. La universidad y el trabajo son colaborativos. Si no aprendes ahora, te cuesta después. (2) \emph{``Mejor que cada uno haga su parte y juntamos al final''} --- Es válido pero limitado: pierdes los beneficios de la coautoría. (3) \emph{``Los compañeros no responden, mejor lo hago todo yo''} --- Falso. Hablar con el grupo y reorganizar es la salida; trabajar de más por miedo a comunicar es el problema. (4) \emph{``En la nube se pierden los archivos''} --- Falso. Es al revés: la nube es la mejor protección contra perder. (5) \emph{``No necesito cuenta institucional, uso mi correo personal''} --- Falso. La cuenta institucional te da Office 365 gratis, espacio en OneDrive y respaldo. Aprovéchala.
\end{softbox}

\begin{infoband}{milcTurquesa}
\textbf{Lo que va al cuaderno (Actividad 2 · EXPLICA).} Título: «Actividad 2 --- 5 principios + ruta P1». \textbf{Formato:} bloque A con los 5 principios numerados y un ejemplo aplicado a digital de cada uno. Bloque B con la lista de 10 sesiones del periodo. \textbf{Extensión:} 15 líneas. \textbf{Sabes que terminaste cuando} podrías explicarle a tu mamá la idea de minga digital en 2 minutos.
\end{infoband}

\puente{Con los principios claros, organizas tu cuaderno y firmas el compromiso.}

\begin{titlebox}{milcMagenta}
Fase 3 · Praxis con pensamiento computacional
\end{titlebox}

\begin{softbox}{milcMagenta}{milcGris}{Construyendo el producto con los cuatro pilares}
\textbf{Actividad 3 · CREA --- Cuaderno inaugural + 5 preguntas-radar + investigación sobre la minga} (30 min · individual). Vas a dejar tu cuaderno listo para todo el periodo 1 y plantear las preguntas que esperas que se respondan.
\end{softbox}

\small
\begin{tabularx}{\textwidth}{>{\bfseries\RaggedRight\arraybackslash}p{3.6cm}Y}
\rowcolor{milcMagenta!90}\color{white}Pilar & \color{white}\bfseries Aplicación al producto\\
Descomposición & \textbf{Paso 1 --- Portada del cuaderno de grado 7.} En la primera hoja escribe en grande: \textbf{``Cuaderno de Tecnología --- Grado 7''}. Nombre completo, grado y letra (7A, 7B, etc.), I.E. Sor María Juliana, año 2026. Decora a tu gusto pero que se lea claro desde 1 metro de distancia.\\
Patrones & \textbf{Paso 2 --- Índice del periodo 1.} En la segunda hoja escribe el título: \textbf{``Índice --- Periodo 1: Trabajo colaborativo en la nube''}. Lista numerada de las 10 sesiones (S1 a S10) con sus títulos. Deja una columna a la derecha vacía para anotar la fecha cuando veas cada sesión.\\
Abstracción & \textbf{Paso 3 --- Compromiso del año (los 5 acuerdos).} En la tercera hoja escribe \textbf{``Mis 5 acuerdos del trabajo colaborativo''}. Reproduce los 5 acuerdos (responsabilidad, puntualidad, respeto, escucha, comunicación clara). Al lado de cada uno, palomita ✓ si te comprometes. Si dudas con alguno, conversa con el profe antes de firmar.\\
Algoritmo de armado & \textbf{Paso 4 --- Mis 5 preguntas-radar.} En la misma hoja, debajo de los acuerdos, escribe \textbf{``Mis 5 preguntas-radar del P1''} y una lista numerada de 5 preguntas tuyas. Sugerencias: \emph{``¿Cómo se comparte un archivo en OneDrive?''}, \emph{``¿Cómo evitamos que 3 personas borren lo que escribieron al tiempo?''}, \emph{``¿Qué hago si alguien del equipo no aporta?''}. Inventa tus propias 5.\\
\end{tabularx}
\normalsize

\begin{drawbox}{milcMagenta}{Antes de cerrar la S1 del periodo 1 de grado 7}
\checkbox La portada tiene nombre, grado, año y materia legibles. \\[1mm]
\checkbox El índice del periodo 1 está con 10 sesiones. \\[1mm]
\checkbox Los 5 acuerdos colaborativos tienen firma. \\[1mm]
\checkbox Tengo 5 preguntas-radar reales mías. \\[1mm]
\checkbox La mini-investigación tiene mínimo 5 datos concretos. \\[1mm]
\checkbox La sesión 1 tiene encabezado y registro de las 3 actividades.
\end{drawbox}

\begin{softbox}{milcMagenta}{milcGris}{Plantilla de trabajo}
\textbf{Estructura.} \textit{Hoja 1:} portada. \textit{Hoja 2:} índice P1 (10 sesiones). \textit{Hoja 3:} 5 acuerdos colaborativos + 5 preguntas-radar. \textit{Hoja 4:} mini-investigación minga (5+ datos). \textit{Cuaderno principal:} sesión 1 con encabezado fijo.
\end{softbox}

\begin{infoband}{milcMagenta}
\textbf{Lo que va al cuaderno (Actividad 3 · CREA).} Título: «Actividad 3 --- Cuaderno inaugural G7-P1». \textbf{Formato:} 4 hojas iniciales + sesión 1. \textbf{Extensión:} 4 hojas + registro de sesión. \textbf{Sabes que terminaste cuando} tu cuaderno está listo para empezar la S2 sin pérdida de tiempo.
\end{infoband}

\puente{El producto es tu cuaderno inaugural + 5 preguntas + mini-investigación.}

\begin{titlebox}{milcMostaza}
Producto de la sesión
\end{titlebox}
\textbf{Cuaderno inaugural G7·P1 + 5 acuerdos + 5 preguntas + investigación minga · firmado}.
\begin{softbox}{milcMostaza}{milcGris}{Criterios de calidad}
(1) Portada con datos legibles. (2) Índice del P1 con 10 sesiones. (3) 5 acuerdos firmados con palomita. (4) 5 preguntas-radar personales. (5) Mini-investigación con 5+ datos sobre la minga (de fuente humana o internet evaluada).
\end{softbox}

\puente{Antes de salir, verifica que tu cuaderno tiene los 5 elementos completos.}

\begin{titlebox}{milcVino}
Fase 4 · Evaluación liberadora — cinco dimensiones
\end{titlebox}

\begin{softbox}{milcVino}{milcGris}{Sentido de la evaluación}
Evaluar no es solo revisar una nota. Es reconocer qué aprendí, qué cambió en mí, qué emociones manejé, cómo me relaciono con la ciudad y la comunidad, y qué le dejaré a quien viene detrás.
\end{softbox}

\textbf{Mi reflexión liberadora en cinco dimensiones.} Completa cada frase iniciada:

\small
\begin{tabularx}{\textwidth}{>{\bfseries\RaggedRight\arraybackslash}p{3.4cm}Y>{\RaggedRight\arraybackslash}p{4.6cm}}
\rowcolor{milcVino}\color{white}Dimensión & \color{white}\bfseries Mi respuesta (frame starter) & \color{white}\bfseries Algo que descubrí\\
1. Personal & Lo que más cambió en mí fue \linead{6.5cm} & \linead{4.2cm}\\
2. Emocional & La emoción más fuerte fue \linead{2.5cm} y la regulé \linead{3cm} & \linead{4.2cm}\\
3. Ciudadana & En democracia esto importa porque \linead{6cm} & \linead{4.2cm}\\
4. Local (Cartago) & En mi barrio sería útil para \linead{6.5cm} & \linead{4.2cm}\\
5. Intergeneracional & A mi abuela/abuelo le diría \linead{6.5cm} & \linead{4.2cm}\\
\end{tabularx}
\normalsize

\begin{infoband}{milcVino}
\textbf{Para la próxima sustentación.} Vuelve a esta tabla en seis meses. Verás que las respuestas cambian — no porque lo aprendido cambie, sino porque tú cambias. La evaluación liberadora es una conversación contigo a través del tiempo.
\end{infoband}


\puente{Cerramos con 3 voces que te ayudan a entender por qué trabajar juntos es habilidad clave de la vida adulta.}

\begin{titlebox}{milcGradeDark}
Triángulo de pensamiento · cierre filosófico
\end{titlebox}

\begin{softbox}{milcVino}{milcGris}{Dussel · lente del nosotros}
\textit{````Solo en comunidad somos plenamente humanos. El que cree que puede todo solo, todavía no ha entendido.''''}

Hay una idea moderna de \textbf{individualismo} que dice: \emph{``si quieres algo, hazlo tú solo''}. Esa idea hace daño. Los grandes proyectos del mundo --- la independencia de Colombia, la construcción de Cartago, las cosechas del café, la pirámide de Egipto --- los hicieron \emph{equipos de personas}, no genios solitarios. Aprender a trabajar en equipo te quita la presión de \emph{``poder con todo''} y te abre el mundo de \emph{``poder con otros''}. La minga lo sabía. Microsoft 365 lo redescubrió. Tú aprendes esa sabiduría antigua con herramientas nuevas.

\textbf{¿Me he creído alguna vez la idea de que ``mejor solo''? ¿De dónde viene esa idea?}
\end{softbox}
\linead{\linewidth}\\[1mm]
\linead{\linewidth}

\begin{softbox}{milcMostaza}{milcGris}{Estoicismo · Marco Aurelio (emperador que coordinaba imperios) · cuidado interior}
\textit{````El que sabe trabajar solo es bueno. El que sabe trabajar con otros es indispensable.''''}

Marco Aurelio dirigía un imperio. \textbf{No podía hacer nada solo}: dependía de gobernadores, generales, escribanos, médicos, traductores. Aprendió a coordinar, a delegar, a comunicar. Esa habilidad lo hizo grande. Tu vida no será imperial, pero \emph{también vas a depender de equipos} --- en universidad, en trabajo, en tu propia familia. Aprender hoy a trabajar con otros es invertir en tu yo de 30 años.

\textbf{¿Estoy entrenándome para ser ``bueno solo'' o ``indispensable con otros''?}
\end{softbox}
\linead{\linewidth}\\[1mm]
\linead{\linewidth}

\begin{softbox}{milcTurquesa}{milcGris}{Floridi · lente de la infoesfera}
\textit{````La nube no es magia: es la versión digital de algo que la humanidad ya hacía hace siglos --- compartir trabajo.''''}

Cuando enciendes Microsoft 365 y entras a OneDrive, no estás haciendo algo radicalmente nuevo. Estás haciendo \textbf{una minga con tus compañeros de equipo}. La herramienta es nueva; la práctica es antigua. Reconocer esa continuidad te da \emph{paz}: no estás solo, ni eres la primera generación que enfrenta esto, ni la herramienta te define. Tú la usas para hacer lo que los humanos llevan haciendo siempre: trabajar juntos.

\textbf{¿Veo la nube como magia que me asusta, o como herramienta que sirve a una práctica antigua?}
\end{softbox}
\linead{\linewidth}\\[1mm]
\linead{\linewidth}

\begin{titlebox}{milcVino}
Compromiso final
\end{titlebox}

\small
\begin{tabularx}{\textwidth}{>{\RaggedRight\arraybackslash}p{3.0cm}YYY}
\rowcolor{milcVino}\color{white}\bfseries Criterio & \color{white}\bfseries Lo logré & \color{white}\bfseries En proceso & \color{white}\bfseries Necesito apoyo\\
Comprensión del tema & Explico ideas centrales con mis palabras. & Reconozco ideas pero falta precisión. & Necesito repasar conceptos base.\\
Pensamiento computacional & Apliqué los 4 pilares al diseñar mi producto. & Apliqué algunos pilares. & Necesito releer la fase 2.\\
Aplicación y producto & Mi producto cumple los criterios y se puede revisar. & Producto funcional con ajustes pendientes. & Aún en construcción.\\
Reflexión 5 dimensiones & Respondí cada dimensión con honestidad. & Respondí algunas con detalle. & Mis respuestas son muy generales.\\
Triángulo de pensamiento & Conecté las 3 voces con mi trabajo real. & Conecté al menos una. & No relaciono las voces con mi proceso.\\
\end{tabularx}
\normalsize

\textbf{Hoy entendí que:}\\[1mm]
\linead{\linewidth}\\[1mm]
\linead{\linewidth}

\textbf{Mi compromiso para la próxima sesión es:}\\[1mm]
\linead{\linewidth}\\[1mm]
\linead{\linewidth}

\begin{softbox}{milcMostaza}{milcGris}{Cita de despedida · Marco Aurelio}
\textit{``El obstáculo en el camino se vuelve el camino. Lo que estorba al camino se vuelve camino.''}\\[1mm]
Cada error de hoy, bien observado, se convierte en pieza del camino que pisarás mañana.
\end{softbox}

\small
\begin{tabularx}{\textwidth}{>{\bfseries}p{3.6cm}Y}
\rowcolor{milcGradeDark!12}Nombre completo: & \linead{10cm}\\
Fecha: & \linead{4cm} \quad Curso/grupo: \linead{4cm}\\
Mi firma: & \linead{9cm}\\
\end{tabularx}
\normalsize

\begin{infoband}{milcGradeDark}
\textbf{Para la próxima sesión.} Termina tu mini-investigación sobre la minga si te quedó pendiente (pregunta en casa). Trae tu cuaderno listo y la curiosidad encendida. La próxima clase es S2: \emph{Microsoft 365 --- qué es la nube}. Vas a conocer las 7-8 aplicaciones del paquete y a entender por qué se llama \emph{``la nube''}. Hoy te recibiste como aprendiz del oficio colaborativo; la próxima clase abrimos el taller digital.
\end{infoband}

\end{document}
